\chapter{动理论与理想气体}\label{ch:b1:kinetic-theory}

教室里的空气，重量与一个成年人相当，所含分子数
比地球上的沙粒数还多；每个分子以步枪子弹的速率飞行，
每秒与邻近分子碰撞一百亿次。
没有人能追踪其中任何一个分子，也没有人需要这样做。一组
平均量——\omterm{def:b1:kinetic-theory:pressure}{压强}、\omterm{def:b1:kinetic-theory:temperature}{温度}、体积、物质的量——
便足以在工程的各种用途中完整描述空气，
分子撞击器壁这一简单图像，既解释了这些平均量从何而来，
也说明了它们之间的关系。本章开启本卷的热力学部分，
从两端来看物质：微观的
混乱与宏观的平静，
以及在\omterm{def:b1:kinetic-theory:model}{理想气体}情形下把二者联系起来的动理论。

\section{两种尺度的描述}

\begin{definition}[微观尺度与宏观尺度]\label{def:b1:kinetic-theory:scales}
\index{微观尺度}\index{宏观尺度}\index{阿伏伽德罗常数}\index{摩尔}
在\emph{微观}尺度上，物质由分子组成：$\qty{1}{mol}$
含有 $N_A = \qty{6.02e23}{}$ 个分子，每个分子的直径约为 \qty{0.3}{nm}；
在常温常压下的气体中，分子间距约为 \qty{3}{nm}，
分子以每秒数百米的速率运动，大约每 \qty{70}{nm} 的路程与其他分子碰撞一次。
在\emph{宏观}尺度上，一个系统由少数几个
\emph{\omterm{def:b1:kinetic-theory:equilibrium}{状态参量}}描述——\omterm{def:b1:kinetic-theory:pressure}{压强} $P$、体积 $V$、\omterm{def:b1:kinetic-theory:temperature}{温度} $T$、
物质的量 $n$——它们都是大量分子的平均结果。
若一个量在系统加倍时也随之加倍，则称其为\emph{广延量}
（$V$、$n$、质量、能量）；若不加倍，则称为\emph{强度量}
（$P$、$T$、密度）。
\end{definition}

\begin{definition}[热力学平衡]\label{def:b1:kinetic-theory:equilibrium}
\index{热力学平衡}\index{状态参量}
系统处于\emph{热力学平衡}，是指其状态参量
均匀且不随时间变化，并且没有宏观的物质流或能量流穿过系统。
只有在此条件下，$P$ 和 $T$ 才对整个系统有确定的意义；
\emph{状态方程}则将它们与 $V$、$n$ 联系起来。
\end{definition}

\begin{definition}[压强]\label{def:b1:kinetic-theory:pressure}
\index{压强}\index{帕斯卡}\index{巴}
静止的流体对器壁的每个面元 $\dd S$ 施加一个力，
$\dd\vect F = P\,\dd S\,\vect n$，该力垂直于面元并指向器壁；$P$ 称为
\emph{压强}，单位为帕斯卡（$\qty{1}{Pa} = \qty{1}{N/m^2}$）；$\qty{1}{bar}
= \qty{1e5}{Pa}$，$\qty{1}{atm} = \qty{1.013e5}{Pa}$。在给定点处，
压强不依赖于面元的取向（见\cref{ch:b1:fluid-statics}）。
\end{definition}

\begin{definition}[温度]\label{def:b1:kinetic-theory:temperature}
\index{温度}\index{开尔文}\index{热力学第零定律}
两个系统通过允许能量通过的器壁接触，经过一段时间后达到
共同的状态：此时它们具有相同的\emph{温度}；
若两个物体各自与第三个物体处于平衡，则它们彼此也处于平衡
（\emph{热力学第零定律}）。开尔文温标由
$k_B = \qty{1.380649e-23}{J/K}$ 精确定义；摄氏温度 $\theta = T - 273.15$。
下文的动理论将赋予 $T$ 微观意义。
\end{definition}

\section{理想气体的动理论模型}

\begin{definition}[理想气体的微观模型]\label{def:b1:kinetic-theory:model}
\index{理想气体}\index{理想气体}
\emph{理想气体}是由 $N$ 个分子组成的集合，
这些分子视为质量为 $m$ 的质点，在体积为 $V$ 的容器中自由且无规则地运动，
除短暂碰撞外没有相互作用。
在平衡态下，速度分布是各向同性且定态的；
$n^* = N/V$ 为数密度，$\langle v^2\rangle$ 为
均方速率。
\end{definition}

\begin{theorem}[动理论压强]\label{thm:b1:kinetic-theory:pressure}
\index{动理论压强}
气体施加在器壁上的\omterm{def:b1:kinetic-theory:pressure}{压强}为
\[
P = \tfrac13\,n^*m\langle v^2\rangle = \tfrac23\,n^*\langle\epsilon_k\rangle ,
\]
其中 $\langle\epsilon_k\rangle = \tfrac12 m\langle v^2\rangle$ 为
一个分子的平均\omterm{thm:b1:work-and-energy:ke}{动能}。
\end{theorem}

\begin{proof}[部分证明]
取简化模型：所有分子的速率均为 $u$，且分别沿三个坐标轴运动，
每个方向上的分子各占六分之一。一个分子
正面撞击器壁 $x = 0$ 时发生弹性反弹，
并给器壁传递动量 $2mu$。在时间 $\dd t$ 内，朝向器壁运动的分子
能够到达面积 $S$ 的，是位于厚度为 $u\,\dd t$ 的薄层内的分子，
即该处 $n^*Su\,\dd t$ 个分子中的六分之一。力 $= $
单位时间的动量变化 $= (n^*Su\,\dd t/6)(2mu)/\dd t = \tfrac13 n^*mu^2S$，因此
$P = \tfrac13 n^*mu^2$。对速率与方向的实际分布取平均，
将 $u^2$ 换为 $\langle v^2\rangle$（各向同性给出
$\langle v_x^2\rangle = \tfrac13\langle v^2\rangle$）；本科二年级分册
给出完整推导。
\end{proof}

\begin{omfigure}
\begin{tikzpicture}[font=\small, x=1cm, y=1cm]
  \draw[thick] (0,-1.6) -- (0,1.6); \fill[pattern=north east lines] (-0.2,-1.6) rectangle (0,1.6);
  \draw[dashed] (1.8,-1.6) -- (1.8,1.6);
  \draw[Stealth-Stealth] (0,-1.85) -- (1.8,-1.85) node[midway, below] {$u\,\dd t$};
  \foreach \p in {(0.6,1.1),(1.2,0.5),(0.9,-0.2),(1.5,-0.9),(0.4,-1.2),(1.3,1.3)} {\fill[omDef] \p circle (2pt);}
  \draw[-Stealth, omDef, thick] (1.2,0.5) -- (0.5,0.5);
  \draw[-Stealth, omThm, thick] (0.25,0.2) -- (0.95,0.2);
  \node[omDef, font=\footnotesize, above] at (0.85,0.55) {$m\vect u$};
  \node[omThm, font=\footnotesize, below] at (0.6,0.15) {$-m\vect u$};
  \node[right, font=\footnotesize] at (2.0,0.9) {每次撞击传给器壁的};
  \node[right, font=\footnotesize] at (2.0,0.55) {动量 $2mu$};
  \node[right, font=\footnotesize] at (2.0,-0.5) {$\tfrac16 n^*Su\,\dd t$ 次撞击（$\dd t$ 内）};
  \begin{scope}[xshift=7.2cm, yshift=-1.6cm]
  \begin{axis}[omaxis, axis x line=bottom, axis y line=left, width=6.4cm, height=4.2cm, xlabel={$v$ (\unit{m/s})}, ylabel={$f(v)$},
      xlabel style={at={(axis description cs:0.5,-0.1)}, anchor=north},
      xmin=0, xmax=1400, ymin=0, ymax=1.2, xtick={500,1000}, ytick=\empty, anchor=origin, at={(0,0)}]
    \addplot[omDef, thick, domain=0:1400, samples=120] {(x/422)^2*exp(1-(x/422)^2)};
    \addplot[omThm, thick, domain=0:1400, samples=120] {0.707*(x/597)^2*exp(1-(x/597)^2)};
    \node[omDef, font=\footnotesize] at (axis cs:430,1.12) {\qty{300}{K}};
    \node[omThm, font=\footnotesize] at (axis cs:900,0.72) {\qty{600}{K}};
  \end{axis}
  \end{scope}
\end{tikzpicture}

{\small 左图：\omterm{def:b1:kinetic-theory:pressure}{压强}的动理学起源——分子撞击
器壁并弹回，将动量传给器壁；单位时间内的碰撞次数
乘以 $2mu$ 即为作用力。右图：实际速率呈分布（麦克斯韦
分布，此处取两种\omterm{def:b1:kinetic-theory:temperature}{温度}下的氮气）；模型中的单一速率
$u$ 应换成\omterm{def:b1:kinetic-theory:kinetictemp}{方均根速率}。}
\end{omfigure}

\begin{definition}[动理学温度；状态方程]\label{def:b1:kinetic-theory:kinetictemp}
\index{动理学温度}\index{状态方程!理想气体}\index{气体常量}\index{方均根速率}
\omterm{def:b1:kinetic-theory:model}{理想气体}的\omterm{def:b1:kinetic-theory:temperature}{温度}由其分子的平均
\omterm{thm:b1:work-and-energy:ke}{动能}\emph{定义}，
\[
\langle\epsilon_k\rangle = \tfrac12 m\langle v^2\rangle = \tfrac32 k_BT ,
\]
于是结合\cref{thm:b1:kinetic-theory:pressure}，
\[
PV = Nk_BT = nRT, \qquad R = N_Ak_B = \qty{8.314}{J/(mol.K)} ,
\]
即\emph{\omterm{def:b1:kinetic-theory:model}{理想气体}状态方程}。\emph{方均根
速率}为 $u = \sqrt{\langle v^2\rangle} = \sqrt{3k_BT/m} = \sqrt{3RT/M}$，
$M$ 为\omterm{def:b1:kinetic-theory:scales}{摩尔}质量。
\end{definition}

\begin{example}[空气中的数值]\label{ex:b1:kinetic-theory:air}
在 \qty{300}{K} 下，氮气：$u = \sqrt{3 \times 8.314 \times 300/0.028} =
\qty{517}{m/s}$；氢气为 $\qty{1930}{m/s}$；质量越大，速率越小，
即与 $1/\sqrt M$ 成正比。在 \qty{1}{bar} 与 \qty{300}{K} 下，$n^* = P/k_BT = \qty{2.4e25}{m^{-3}}$：
每\omterm{def:b1:kinetic-theory:scales}{摩尔} \qty{24}{L}，每升中有 \qty{2.7e22}{} 个分子。一间
\qty{200}{m^3} 的教室：\qty{8300}{mol}，\qty{240}{kg} 空气。
\end{example}

\begin{remark}[模型为何成立，及何时失效]\label{rem:b1:kinetic-theory:validity}
\index{平均自由程}
两次碰撞之间，空气分子飞过一段\emph{\omterm{rem:b1:kinetic-theory:validity}{平均自由程}}
$\lambda \approx 1/(\sqrt2\pi d^2n^*) \approx \qty{70}{nm}$——约为
分子大小的两百倍：气体大部分为空，相互作用稀少，
且 $PV = nRT$ 在常压下精确到百分之一以内成立。在高压
或低温下，分子自身的体积及其间的
吸引力不可忽略：范德瓦耳斯方程 $(P + an^2/V^2)(V - nb) = nRT$
同时修正了这两方面；低于临界\omterm{def:b1:kinetic-theory:temperature}{温度}时，气体便会液化
（见\cref{ch:b1:phase-changes}）。
\end{remark}

\begin{omfigure}
\begin{tikzpicture}
\begin{axis}[omaxis, axis x line=bottom, axis y line=left, width=10cm, height=6cm, xlabel={$V$}, ylabel={$P$},
    xlabel style={at={(axis description cs:0.5,-0.08)}, anchor=north},
    xmin=0, xmax=5.2, ymin=-0.4, ymax=4.2, xtick=\empty, ytick=\empty]
  \addplot[omDef, thick, domain=0.3:5.2, samples=100] {1.0/x};
  \addplot[omDef, thick, domain=0.45:5.2, samples=100] {1.5/x};
  \addplot[omDef, thick, domain=0.6:5.2, samples=100] {2.0/x};
  \addplot[omThm, thick, domain=0.75:5.2, samples=100] {2.5/x};
  \node[omDef, font=\footnotesize, below] at (axis cs:4.9,0.19) {$T_1$};
  \node[omThm, font=\footnotesize, above] at (axis cs:4.9,0.53) {$T_4$};
  \draw[-Stealth, omExo] (axis cs:1.3,1.0) -- (axis cs:2.3,2.0) node[above right, font=\footnotesize] {$T$ 增大};
  \node[font=\footnotesize] at (axis cs:3.4,3.2) {$PV = nRT$：等温线 $T_1 < \dots < T_4$};
\end{axis}
\end{tikzpicture}

{\small \omterm{def:b1:kinetic-theory:model}{理想气体}的克拉珀龙图：每一\omterm{def:b1:kinetic-theory:temperature}{温度}下的
等温线是双曲线 $P = nRT/V$，$T$ 越高则越靠上。一个过程
是该平面中的一条路径；循环即闭合回路（见\cref{ch:b1:heat-engines}）。}
\end{omfigure}

\section{理想气体的内能}

\begin{definition}[内能]\label{def:b1:kinetic-theory:U}
\index{内能}
系统的\emph{内能} $U$ 为其分子在系统静止的
\omterm{def:b1:point-kinematics:frame}{参考系}中的\omterm{thm:b1:work-and-energy:ke}{动能}，与其分子间\omterm{def:b1:work-and-energy:conservative}{势能}之和；
它是一个广延状态函数。
\end{definition}

\begin{proposition}[理想气体的内能]\label{prop:b1:kinetic-theory:Ugas}
\index{单原子气体}\index{双原子气体}\index{定容热容}\index{焦耳定律}
对\omterm{def:b1:kinetic-theory:model}{理想气体}，$U$ 仅依赖于 $T$（第一焦耳定律）：
\begin{gather*}
U = \tfrac32 nRT \quad (\text{单原子：He, Ne, Ar}), \\
U = \tfrac52 nRT \quad (\text{常温下的双原子：N}_2, \text{O}_2, \text{H}_2),
\end{gather*}
因此定容\omterm{def:b1:kinetic-theory:scales}{摩尔}热容 $C_{V,m} = \dd U_m/\dd T$
分别为 $\tfrac32 R = \qty{12.5}{J/(mol.K)}$ 与
$\tfrac52 R = \qty{20.8}{J/(mol.K)}$。
\end{proposition}

\begin{proof}[部分证明]
依假设，分子间无相互作用\omterm{def:b1:work-and-energy:conservative}{势能}，故 $U = N\langle\epsilon_k\rangle
= \tfrac32 Nk_BT$，对质点分子成立。双原子分子还会转动：
平衡时，其两个转动自由度各携带
与每个平动自由度相同的 $\tfrac12 k_BT$（\emph{能量均分}
定理，在本阶段不加证明地承认；本科三年级分册推导了它，并解释
为何振动仅在一千\omterm{def:b1:kinetic-theory:temperature}{开尔文}以上才参与）。
\end{proof}

\begin{example}[一间教室的能量]\label{ex:b1:kinetic-theory:room}
\qty{300}{K} 下教室中的 \qty{8300}{mol} 空气：$U = \tfrac52 \times
8300 \times 8.314 \times 300 = \qty{5.2e7}{J}$——相当于一升半汽油的能量，
这些能量以分子运动的形式储存。将教室\omterm{def:b1:kinetic-theory:temperature}{温度}升高 $1$ K
需 $\tfrac52 nR = \qty{170}{kJ}$（定容过程）。
\end{example}

\section{凝聚相}

\begin{proposition}[不可压缩、不可膨胀相的模型]\label{prop:b1:kinetic-theory:condensed}
\index{凝聚相}\index{热膨胀}\index{压缩系数}
在本阶段，液体与固体被视为体积固定的相——
既不可压缩、也不可膨胀——其\omterm{def:b1:kinetic-theory:U}{内能}
仅取决于 $T$：$\dd U = C\,\dd T$，其中 $C = mc$，$c$ 为
比热容（水为 \qty{4180}{J/(kg.K)}，钢为
\qty{450}{J/(kg.K)}）。实际中微小的偏差由\omterm{prop:b1:kinetic-theory:condensed}{热膨胀}
系数 $\alpha = (1/V)(\partial V/\partial T)_P$（金属中为
$\sim\qty{1e-5}{K^{-1}}$，液体中为 \qty{2e-4}{K^{-1}}）；以及等温
\omterm{prop:b1:kinetic-theory:condensed}{压缩系数} $\chi_T = -(1/V)(\partial V/\partial P)_T$（水中 $\sim\qty{5e-10}{Pa^{-1}}$，
与气体的 $1/P \approx \qty{1e-5}{Pa^{-1}}$ 相比）。
\end{proposition}
\admitted

\begin{example}[水有多难压缩]\label{ex:b1:kinetic-theory:water}
马里亚纳海沟底部的水（\qty{1100}{bar}）被压缩了
约 $\chi_T\Delta P \approx 5\times10^{-10} \times 1.1\times10^8 = 5\%$；一根
\qty{30}{m} 长的钢轨升温 $50$ K，其伸长量为 $\alpha L\Delta T =
\qty{18}{mm}$——这正是设置伸缩缝的原因。与气体相比，
其体积在\omterm{def:b1:kinetic-theory:pressure}{压强}加倍时减半，因此固定体积的模型
在这一体积范围的大多数情形中，误差仅为百分之几。
\end{example}

\section{练习}

\begin{exercise}[$\star$]\label{exo:b1:kinetic-theory:1}
求 \qty{1.0}{L} 气体在 \qty{0}{\degreeCelsius} 和
\qty{1}{atm} 下的分子数；求相邻分子间的平均距离，
并与分子大小 \qty{0.3}{nm} 作比较。
\end{exercise}

\begin{exercise}[$\star$]\label{exo:b1:kinetic-theory:2}
求 H$_2$、N$_2$、O$_2$ 和 CO$_2$ 在 \qty{300}{K} 时的\omterm{def:b1:kinetic-theory:kinetictemp}{方均根速率}，
以及 N$_2$ 在 \qty{77}{K}（液氮的沸点）时的\omterm{def:b1:kinetic-theory:kinetictemp}{方均根速率}。
\end{exercise}

\begin{exercise}[$\star$]\label{exo:b1:kinetic-theory:3}
一只体积为 \qty{30}{L} 的汽车轮胎，在 \qty{20}{\degreeCelsius} 下充气至
表压 \qty{2.5}{bar}。求其内空气的物质的量与质量；
行驶一段路程、轮胎被加热至 \qty{50}{\degreeCelsius} 后，求其绝对\omterm{def:b1:kinetic-theory:pressure}{压强}。
\end{exercise}

\begin{exercise}[$\star$]\label{exo:b1:kinetic-theory:4}
估算空气分子在 \qty{1}{bar}、\qty{300}{K} 下的\omterm{rem:b1:kinetic-theory:validity}{平均自由程}，
（$d = \qty{0.37}{nm}$），以及单个分子的碰撞频率。
在何种\omterm{def:b1:kinetic-theory:pressure}{压强}下，\omterm{rem:b1:kinetic-theory:validity}{平均自由程}达到 \qty{1}{m}（即“真空”）？
\end{exercise}

\begin{exercise}[$\star\star$]\label{exo:b1:kinetic-theory:5}
求 \qty{1.013}{bar}、\qty{20}{\degreeCelsius} 下空气的密度
（$M = \qty{29}{g/mol}$），再求 \qty{100}{\degreeCelsius} 下的密度；
并求体积为 \qty{2800}{m^3} 的热气球的升力（即两种\omterm{def:b1:kinetic-theory:temperature}{温度}下空气的重量之差）。
\end{exercise}

\begin{exercise}[$\star\star$]\label{exo:b1:kinetic-theory:6}
空气中按分子数计 $21\%$ 为 O$_2$，求海平面处氧分压；
求海拔 \qty{5000}{m}（$P = \qty{0.54}{atm}$）处的氧分压，
以及珠穆朗玛峰峰顶（\qty{0.33}{atm}）处的氧分压。为什么人体在那里会感到吃力？
\end{exercise}

\begin{exercise}[$\star\star$]\label{exo:b1:kinetic-theory:7}
在 \qty{300}{K} 下，将 1 \omterm{def:b1:kinetic-theory:scales}{摩尔} CO$_2$ 置于 \qty{1.0}{L} 的容器中：
先由\omterm{def:b1:kinetic-theory:model}{理想气体}状态方程求\omterm{def:b1:kinetic-theory:pressure}{压强}，再由范德瓦耳斯方程求\omterm{def:b1:kinetic-theory:pressure}{压强}（$a = \qty{0.364}{Pa.m^6/mol^2}$，
$b = \qty{4.27e-5}{m^3/mol}$）。在此情形中，哪一项修正占主导？
\end{exercise}

\begin{exercise}[$\star\star$]\label{exo:b1:kinetic-theory:8}
求体积为 \qty{50}{m^3} 的房间内空气在 \qty{1}{bar}、
\qty{300}{K} 下的\omterm{def:b1:kinetic-theory:U}{内能}；求定容升温 \qty{5}{K} 所需的能量；
并求用 \qty{2}{kW} 加热器（无能量损失）加热所需的时间。
\end{exercise}

\begin{exercise}[$\star\star$]\label{exo:b1:kinetic-theory:9}
求一根 \qty{30}{m} 长的钢轨（$\alpha = \qty{1.2e-5}{K^{-1}}$）
在 \qty{-10}{\degreeCelsius} 至 \qty{40}{\degreeCelsius} 间的长度变化；
一个黄铜环（$\alpha = \qty{1.9e-5}{K^{-1}}$），其内径为
\qty{49.95}{mm}，需套在直径为 \qty{50.00}{mm} 的轴上，
求它至少需要升高多少\omterm{def:b1:kinetic-theory:temperature}{温度}？
\end{exercise}

\begin{exercise}[$\star\star\star$]\label{exo:b1:kinetic-theory:10}
求 \qty{300}{K} 时每个分子的平均能量（$\tfrac32 k_BT$），分别以
焦耳与\omterm{def:b1:charged-particles:eV}{电子伏特}表示；求氩和氮的\omterm{def:b1:kinetic-theory:scales}{摩尔}热容 $C_{V,m}$；为什么氮的
\omterm{def:b1:kinetic-theory:scales}{摩尔}热容更大，尽管其分子更轻？
\end{exercise}

\begin{exercise}[$\star\star\star$]\label{exo:b1:kinetic-theory:11}
在六方向模型中重新推导\omterm{def:b1:kinetic-theory:pressure}{压强}的动理学表达式，
然后证明：对于各向同性的速度分布，$\langle v_x^2\rangle
= \tfrac13\langle v^2\rangle$，且\omterm{def:b1:kinetic-theory:pressure}{压强}公式保持不变。
\end{exercise}

\begin{exercise}[$\star\star\star$]\label{exo:b1:kinetic-theory:12}
求 H$_2$ 与 O$_2$ 在 \qty{300}{K} 时的\omterm{def:b1:kinetic-theory:kinetictemp}{方均根速率}，并与
地球的\omterm{prop:b1:central-forces:escape}{逃逸速度}、月球的\omterm{prop:b1:central-forces:escape}{逃逸速度}（\qty{2.4}{km/s}）作比较。
已知速率大于 $3u$ 的分子所占比例很小，但不为
零，解释为什么地球失去了氢，而月球失去了整个
大气层。
\end{exercise}

\section{问题：房间里的空气}

\begin{problem}[{周末问题——一间教室里的空气：统计其分子数、称量其质量，
测量分子的飞行时间，并由同一模型
推算气球的升力、山区空气的稀薄程度与
声速}]
\label{pb:b1:kinetic-theory:1}
一间教室的尺寸为 $10 \times 8 \times \qty{2.5}{m}$；其中空气所处的状态为
$P = \qty{1.013e5}{Pa}$，$T = \qty{293}{K}$，\omterm{def:b1:kinetic-theory:scales}{摩尔}质量 $M = \qty{29.0}{g/mol}$，
按分子数计，$79\%$ 为 N$_2$，$21\%$ 为 O$_2$。$R = \qty{8.314}{J/(mol.K)}$，
$k_B = \qty{1.38e-23}{J/K}$，$N_A = \num{6.02e23}$，$g = \qty{9.81}{m/s^2}$。

\medskip
\noindent\textbf{第一部分——计数与称重。}
\begin{enumerate}
  \item 房间内空气的物质的量（\omterm{def:b1:kinetic-theory:scales}{摩尔}数）及分子数。
  \item 求空气的质量，并与一个人的质量作比较。
  \item 求数密度 $n^*$ 及分子间的平均距离
        （$n^{*-1/3}$），并与分子直径 \qty{0.37}{nm} 作比较。
  \item 求 N$_2$ 与 O$_2$ 的\omterm{def:b1:kinetic-theory:kinetictemp}{方均根速率}；为什么在相同\omterm{def:b1:kinetic-theory:temperature}{温度}下
        较重的气体运动得较慢？
  \item 求一个分子的平均\omterm{thm:b1:work-and-energy:ke}{动能}，以及房间内空气的总平动
        \omterm{thm:b1:work-and-energy:ke}{动能}，并与一辆以
        \qty{100}{km/h} 行驶的汽车的\omterm{thm:b1:work-and-energy:ke}{动能}作比较。
  \item 求\omterm{rem:b1:kinetic-theory:validity}{平均自由程} $\lambda \approx 1/(\sqrt2\pi d^2n^*)$，以及
        一个分子两次碰撞之间的平均时间。
\end{enumerate}

\medskip
\noindent\textbf{第二部分——碰撞产生的\omterm{def:b1:kinetic-theory:pressure}{压强}。}
\begin{enumerate}[resume]
  \item 用六方向模型，取 $u = \qty{500}{m/s}$，计算
        每秒打在 \qty{1}{cm^2} 墙面上的碰撞次数。
  \item 计算每秒传递的动量，并检验它
        是否再现大气\omterm{def:b1:kinetic-theory:pressure}{压强}。
  \item 求空气作用于一面 \qty{8}{m} × \qty{2.5}{m} 墙面上的总力；
        墙为何不移动？
  \item 房间关闭门窗并密封后加热至
        \qty{303}{K}：求新的\omterm{def:b1:kinetic-theory:pressure}{压强}，以及此时墙面所受的合力。
  \item 若开窗并以同样方式加热：离开的是什么，离开多少？
\end{enumerate}

\medskip
\noindent\textbf{第三部分——浮力。}
\begin{enumerate}[resume]
  \item 求房间内空气的密度，以及在相同
        \omterm{def:b1:kinetic-theory:pressure}{压强}下 \qty{373}{K} 时空气的密度。
  \item 一只体积为 \qty{2800}{m^3} 的热气球充满\omterm{def:b1:kinetic-theory:temperature}{温度}为
        \qty{373}{K} 的空气：求热空气的重力、被排开的冷空气的重力，
        以及可用升力（阿基米德，
        \cref{ch:b1:fluid-statics}）。
  \item 气囊、吊篮、燃烧器和燃气的总质量为 \qty{300}{kg}：
        它能吊起多少名质量为 \qty{75}{kg} 的乘客？
  \item 空气应加热到多少\omterm{def:b1:kinetic-theory:temperature}{温度}，才能使升力加倍，
        实际中为何不这样做（气囊由尼龙制成）？
  \item 同体积的氦气球，\omterm{def:b1:kinetic-theory:temperature}{温度}为 \qty{293}{K}
        （$M = \qty{4}{g/mol}$）：能否升空？
        游览用的气球为何仍用热空气？
\end{enumerate}

\medskip
\noindent\textbf{第四部分——稀薄空气与声速。}
\begin{enumerate}[resume]
  \item 求房间内氧气的分\omterm{def:b1:kinetic-theory:pressure}{压强}。
  \item 在珠穆朗玛峰顶，$P = \qty{0.33}{atm}$，$T \approx \qty{250}{K}$：
        求氧分子的数密度，并与房间内的数密度相比较。
  \item 大气\omterm{def:b1:kinetic-theory:pressure}{压强}随高度近似地按
        $P(z) = P_0\eu^{-z/H}$ 规律下降，其中 $H = RT/Mg$（见\cref{ch:b1:fluid-statics}）：
        求 \qty{260}{K} 时的 $H$，以及使 $P = P_0/2$ 的高度。
  \item \omterm{def:b1:kinetic-theory:model}{理想气体}中的声速为 $c_s = \sqrt{\gamma RT/M}$，
        其中空气的 $\gamma = 1.4$（见\cref{ch:b1:first-law}）：求
        \qty{293}{K} 时的声速，并与\omterm{def:b1:kinetic-theory:kinetictemp}{方均根速率}比较；解释二者
        接近的意义。
  \item $c_s$ 如何随 $T$ 变化？求房间内\omterm{def:b1:kinetic-theory:temperature}{温度}为
        \qty{303}{K} 时的声速，以及 \qty{293}{K} 时氦气中的声速（$\gamma = 5/3$）：
        为何吸入氦气后的声音听起来较高？
  \item 在珠穆朗玛峰顶，声速比房间内的
        快还是慢？相差多少？
  \item 房间内的空气在相同的 $P$ 与 $T$ 下被氩气替代：
        分子数、质量和速率各有何变化，
        \omterm{def:b1:kinetic-theory:U}{内能}又如何变化？
  \item 求房间内空气的\omterm{def:b1:kinetic-theory:U}{内能}（$U = \tfrac52 nRT$），
        以及在定容下将其加热 \qty{10}{K} 所需的能量。
  \item 用三行概括动理论模型的成果：
        它解释了哪些宏观量，
        并由哪一个单一的微观平均量导出。
\end{enumerate}
\end{problem}
